\documentclass[11pt,a4paper]{article}
\usepackage[T1]{fontenc}
\usepackage[utf8]{inputenc}
\usepackage{lmodern}
\usepackage[a4paper,margin=27mm,headheight=15pt]{geometry}
\usepackage{amsmath,amssymb,amsthm,mathtools}
\usepackage{microtype}
\usepackage{booktabs,longtable,array,enumitem}
\usepackage{xcolor}
\usepackage{fancyhdr}
\usepackage{hyperref}
\usepackage[nameinlink,noabbrev]{cleveref}
\definecolor{ink}{HTML}{19384C}
\definecolor{accent}{HTML}{286575}
\hypersetup{colorlinks=true,linkcolor=ink,citecolor=accent,urlcolor=accent,
  pdftitle={Surreal and Surcomplex Numbers in Theoretical Physics},
  pdfauthor={Research synthesis prepared with AI assistance},
  pdfsubject={Black-hole singularities, non-Archimedean analysis, and physical interpretation}}
\pagestyle{fancy}
\fancyhf{}
\fancyhead[L]{\small\textsc{Surreal scalars and physical singularities}}
\fancyhead[R]{\small\thepage}
\renewcommand{\headrulewidth}{0.3pt}
\setlength{\parskip}{3.5pt}
\setlength{\parindent}{1.2em}
\setlength{\emergencystretch}{2em}
\setcounter{tocdepth}{2}
\numberwithin{equation}{section}
\newtheorem{theorem}{Theorem}[section]
\newtheorem{proposition}[theorem]{Proposition}
\newtheorem{corollary}[theorem]{Corollary}
\newtheorem{lemma}[theorem]{Lemma}
\theoremstyle{definition}
\newtheorem{definition}[theorem]{Definition}
\newtheorem{example}[theorem]{Example}
\theoremstyle{remark}
\newtheorem{remark}[theorem]{Remark}
\newcommand{\R}{\mathbb R}
\newcommand{\C}{\mathbb C}
\newcommand{\Q}{\mathbb Q}
\newcommand{\Z}{\mathbb Z}
\newcommand{\N}{\mathbb N}
\newcommand{\No}{\mathbf{No}}
\newcommand{\Sc}{\No[i]}
\newcommand{\F}{F}
\newcommand{\K}{K}
\newcommand{\OO}{\mathcal O}
\newcommand{\mm}{\mathfrak m}
\newcommand{\st}{\operatorname{st}}
\newcommand{\supp}{\operatorname{supp}}
\newcommand{\val}{\operatorname{v}}
\newcommand{\Oz}{\mathbf{Oz}}
\newcommand{\Ric}{\operatorname{Ric}}
\newcommand{\Scal}{\operatorname{Scal}}
\newcommand{\Kret}{\mathcal K}
\newcommand{\dd}{\mathrm d}
\newcommand{\ee}{\mathrm e}
\newcommand{\ii}{\mathrm i}
\newcommand{\Id}{\operatorname{Id}}
\newcommand{\ExpEK}{\operatorname{Exp}_{\mathrm{EK}}}
\newcommand{\sinEK}{\operatorname{Sin}_{\mathrm{EK}}}
\newcommand{\cosEK}{\operatorname{Cos}_{\mathrm{EK}}}
\newcommand{\arxiv}[1]{\href{https://arxiv.org/abs/#1}{arXiv:#1}}
\newcommand{\doi}[1]{\href{https://doi.org/#1}{doi:#1}}
\newcommand{\repo}[1]{\href{https://github.com/VladimirReshetnikov/Surreal/blob/39f2be6667ade51bca2b45daa47e289d69c09764/#1}{\nolinkurl{#1}}}

\begin{document}
\begin{titlepage}
\centering
\vspace*{16mm}
{\color{accent}\large\textsc{Mathematical physics / critical research assessment}\par}
\vspace{13mm}
{\color{ink}\Huge\bfseries Surreal and Surcomplex Numbers\par}
\vspace{4mm}
{\color{ink}\LARGE in Theoretical Physics\par}
\vspace{10mm}
{\Large Black-hole singularities, asymptotic structure,\par
and the limits of changing the scalar field\par}
\vspace{13mm}
{\large 21 September 2026\par}
\vspace{13mm}
\begin{minipage}{0.89\textwidth}
\normalsize
\textbf{Central assessment.} Surreal scalars can make hierarchies of divergent and
infinitesimal quantities exact algebraic objects. This is potentially valuable for
singular perturbation theory and selected differential equations. It does not,
by itself, produce a regular black-hole interior, extend an incomplete geodesic,
or supply a quantum theory of gravity. The most defensible near-term approach
keeps an ordinary spacetime and equips its fields with controlled Hahn or
transseries coefficients.

\medskip
\textbf{Scope.} A self-contained mathematical analysis, an audit of the relevant
parts of the \texttt{VladimirReshetnikov/Surreal} repository, worked relativistic
examples, explicit compatibility and obstruction results, and a proposed research
program. Established results, deductions made here, and speculative physical
interpretations are distinguished throughout.
\end{minipage}
\vfill
{\small Prepared with AI assistance. Not a refereed article or a formal verification.\par
Symbolic checks of selected identities accompany the source.\par}
\end{titlepage}

\begin{abstract}
Could Conway's surreal field $\No$, or its algebraic closure $\No[i]$, replace
$\R$ and $\C$ in theoretical physics, especially near a black-hole singularity?
The answer depends on what ``replace'' and ``resolve'' mean. Finite-dimensional
algebra, Lorentzian quadratic forms, and quantum-state algebra extend readily to
suitable real-closed coefficient fields. Differential equations, integration,
causal geometry, and quantum field theory require substantially more structure.
We explain the established surreal differential and integration theories without
conflating their distinct derivatives, summation rules, or domains of validity.

The Schwarzschild example is developed in invariant and proper-time variables.
At an infinitesimal radius its curvature is an exactly representable infinite
number, but its pole, finite infall time, and divergent tidal deformation remain.
A pole-persistence proposition makes the limited nature of scalar extension
precise. A formal reduction theorem shows that an infinitesimal deformation of a
nondegenerate metric has the expected classical Einstein tensor as its zeroth
coefficient. Conversely, a boundary-layer example shows why pointwise standard
part alone cannot justify differentiating a limiting metric. A Hayward-type core
provides an explicit matched-asymptotic test: Hahn summability correctly selects
inner and outer regimes, but does not select or physically justify the modified
metric. We also prove finite-dimensional quantum shadow results and exhibit an
incompatibility between the canonical integer-part global phase and the
Berarducci--Mantova scalar derivation. These results support surreal methods as
a rigorous asymptotic language, not a demonstrated singularity cure or a
validated fundamental scalar replacement.
\end{abstract}

\tableofcontents
\clearpage

\section{The question that must actually be answered}\label{sec:question}

\subsection{Four different projects}
Replacing the real numbers can mean four importantly different things. An
\emph{algebraic extension} lets a metric, a matrix, or a curvature expression take
values in a larger field while keeping the same formulas. An \emph{asymptotic
extension} interprets a small parameter as an infinitesimal and encodes a family
of approximations in a Hahn series or a transseries. A \emph{nonstandard
reformulation} enlarges a structured mathematical universe, retaining specified
transfer and internality principles and recovering ordinary quantities by a
standard-part construction. Finally, a \emph{new physical theory} changes the
allowed states, dynamics, geometry, or observational rules.

These projects are not equivalent. A useful reformulation may make proofs or
computations easier without changing a single prediction. An algebraic
extension can make an expression meaningful at a new nonzero argument without
making it meaningful at its old singular argument. A genuinely new theory may
use surreal numbers, but its physical content comes from the whole theory, not
from the name of its field.

This distinction is especially important for black holes. In classical general
relativity, a singularity theorem generally concludes causal geodesic
incompleteness under geometric and energy assumptions. It does not merely
announce that a computer has encountered a large number, and it need not identify
a curvature-divergent point already belonging to the spacetime
\cite{Penrose,Witten}. Describing the problem as a failure of real arithmetic
misidentifies what needs repair.

\subsection{An operational hierarchy of success}
It is useful to distinguish five increasingly demanding achievements:
\begin{enumerate}[label=\textnormal{(\arabic*)},leftmargin=2.1em,itemsep=2pt]
\item A divergent quantity has a well-defined generalized numerical value.
\item An asymptotic solution has controlled algebraic and differential meaning.
\item A specified class of observers has well-defined observables through a
previously singular regime.
\item The spacetime or its replacement has a suitable global evolution and
completeness or continuation theorem.
\item The model recovers tested physics and makes a discriminating, stable,
empirically meaningful prediction.
\end{enumerate}
Surreal arithmetic can directly assist the first two. The remaining steps need
additional mathematics and physics. Failure to keep these levels separate is the
main source of exaggerated claims about singularity resolution.

\subsection{How strong is the affirmative case?}
There is substantial mathematics behind the affirmative case. Berarducci and
Mantova constructed compatible differential structures on surreal numbers and
related surreal numbers to transseries \cite{BM}. Costin and Ehrlich developed
extension and integration operators for a substantial, explicitly delimited
class of resurgent functions \cite{CE}. These are serious tools for asymptotic
analysis, not just informal manipulations of infinity. Their availability makes
it reasonable to investigate reduced gravitational equations and multiscale
models using surreal methods.

The more ambitious statement---that physical spacetime or quantum amplitudes
fundamentally use $\No$ or $\No[i]$ rather than $\R$ or $\C$---does not follow.
No predictive black-hole completion is established by the repository or the
specific surreal-physics proposal examined here. This is an assessment of the
sources inspected, not a theorem excluding all future surreal-valued theories.

\section{What the repository provides, and what it does not}\label{sec:repo}

The repository was inspected at the pinned snapshot
\begin{center}
\small\texttt{39f2be6667ade51bca2b45daa47e289d69c09764}.
\end{center}
Its documentation catalogue contains fifteen research packages: six on surreal
numbers, seven on surcomplex mathematics, and two on foundations and computation.
The catalogue identifies the reports as AI-assisted drafts. The root README,
separately, describes a growing Lean formalization of generic prerequisites
\cite{RepoRoot,RepoDocs}.

The distinction between the reports and the formalized prerequisites matters.
The described Lean modules cover finite algebra, complexification, polynomial
jets, Hahn summability and standard part, among other results. The root README
explicitly says that this layer does not yet construct the surreal field or
establish the normal-form bridge to Hahn series. It would therefore be wrong
both to call the entire repository unformalized and to treat its successful
Lean build as verification of all its analytic reports. This investigation did
not run the Lean build or audit every declaration.

For the present task, three features are particularly relevant.
First, the analysis report distinguishes Hahn summation from ordinary
convergence and distinguishes several analytic function classes. Second, the
trigonometry report discusses both finite-angle phases and a global
Ehrlich--Kaplan normalization; it does not select a scalar derivation. Third,
the documentation recommends set-sized coefficient workspaces rather than
unrestricted proper-class constructions \cite{RepoAnalysis,RepoTrig,RepoDocs}.

None of these ingredients is already a relativistic spacetime, an Einstein
initial-value theory, a quantum state space with infinite-dimensional spectral
theory, or a rule for measuring surreal-valued observables. A suitable extension
of the repository would therefore add a \emph{physical semantics and a
differential-geometric layer}, not merely more arithmetic identities.

\paragraph{Use made of the repository.}
The analysis below adopts its most useful methodological discipline: every
calculation specifies the coefficient system and the interpretation of a series.
The proofs supplied here do not rely on treating its unrefereed analytic claims
as established physics. The file-level audit is recorded in
\cref{app:audit}.

\section{The scalar systems and their observational shadows}\label{sec:scalars}

\subsection{Surreal numbers are a field, not a collection of error symbols}
Conway's surreal numbers form a proper-class ordered real-closed field containing
$\R$ and the ordinals. Every individual surreal number is a legitimate object
with set-sized normal form, although the collection of all such numbers is not
a set. The algebraic closure obtained by adjoining $i^2=-1$ is
\[
\Sc=\{a+ib:a,b\in\No\}.
\]
These algebraic facts, the normal-form description, and the exponential
structure are discussed in \cite{BM,EK,Ehrlich}.

An infinite positive element such as $\omega$ satisfies $\omega>n$ for every
ordinary natural number $n$. The number $\varepsilon=\omega^{-1}$ is positive
and smaller than every positive real number. Nevertheless,
\[
\varepsilon\ne0,\qquad \varepsilon^{-1}=\omega,
\qquad 0^{-1}\text{ does not exist}.
\]
No field extension can change the last statement while retaining the field
axioms: if $0x=1$, then $0=1$. Likewise, there is no largest positive surreal and
no smallest positive infinitesimal, since $x+1>x$ and
$0<\varepsilon/2<\varepsilon$.

An infinitesimal is therefore not automatically a minimal length, a shortest
time step, or a ultraviolet cutoff. Its presence adds finer scales rather than
removing them. Surreal numbers are also commutative and have no nonzero
nilpotents. Anticommuting Grassmann variables, noncommuting observables, and
operator algebras still require their own constructions.

\subsection{Set-sized Hahn workspaces}
For concrete calculations choose a divisible ordered abelian group $\Gamma$ that
is a set, and use
\begin{equation}\label{eq:hahn}
F_\Gamma=\R((t^\Gamma)),\qquad
K_\Gamma=\C((t^\Gamma))=F_\Gamma[i].
\end{equation}
An element is a formal sum
\[
a=\sum_{\gamma\in S}a_\gamma t^\gamma,
\qquad S\subseteq\Gamma\text{ well ordered}.
\]
The convention is that $t^\gamma$ is small for $\gamma>0$; the valuation of a
nonzero series is $\val(a)=\min\supp(a)$. The ordered field $F_\Gamma$ is
real closed and $K_\Gamma$ is algebraically closed. If $\Gamma$ is realized as
an additive subgroup of $\No$, the normal-form realization sends
$t^\gamma$ to $\omega^{-\gamma}$, consistently with this convention
\cite{BM,EK}.

A family $(a_j)_{j\in J}$ is \emph{strongly Hahn summable} when the union of its
supports is well ordered and each exponent occurs in the support of only
finitely many members. Addition is then coefficientwise. In particular,
\begin{equation}\label{eq:geom}
\sum_{n\ge0}\varepsilon^n=(1-\varepsilon)^{-1}
\end{equation}
is a strong-summation identity for an infinitesimal $\varepsilon$ of positive
valuation. The ordinary convergent series $\sum_{n\ge0}2^{-n}$, viewed as a
family of constant Hahn series, is \emph{not} strongly Hahn summable: exponent
zero occurs infinitely often. Its ordinary real sum can of course be computed
first and embedded afterwards.

Every set of surreal coefficients uses only a set of normal-form exponents.
Taking the divisible additive subgroup generated by those exponents produces a
set-sized workspace containing that set. This observation does not imply that
the workspace is closed under a chosen exponential, a scalar derivation, or
composition; additional closure requirements must be checked. It also does not
make a proper-class indexed field theory into an ordinary set-sized one merely
by changing notation.

\subsection{Finite elements, infinitesimals, and standard part}
Let $F$ be any ordered field extension of $\R$, and write
\[
\OO_F=\{x\in F: |x|<n\text{ for some }n\in\N\},\qquad
\mm_F=\{x\in F: |x|<1/n\text{ for every }n\ge1\}.
\]
These words describe size relative to the embedded reals. An infinite surreal
is an element of a field but is not \emph{finite} in this sense.

\begin{proposition}[Standard part]\label{prop:st}
Each $x\in\OO_F$ has a unique decomposition $x=r+\delta$ with $r\in\R$ and
$\delta\in\mm_F$. The map $\st(x)=r$ is an order-preserving ring homomorphism
with kernel $\mm_F$, inducing $\OO_F/\mm_F\cong\R$.
For $K=F[i]$ it extends componentwise to a homomorphism
$\OO_F+i\OO_F\to\C$ commuting with conjugation.
\end{proposition}
\begin{proof}
The set $A=\{a\in\R:a\le x\}$ is nonempty and bounded above. Let
$r=\sup_\R A$. For every positive real $h$, the defining properties of the
supremum give $r-h<x<r+h$, with harmless weak inequalities at an endpoint.
Thus $x-r$ is infinitesimal. Two distinct real numbers cannot differ by an
infinitesimal, proving uniqueness. Sums of infinitesimals and products of an
infinitesimal with a finite element are infinitesimal. Expanding
$(r+\delta)(s+\eta)$ now proves the homomorphism assertion. Order preservation
and the complex assertion follow directly.
\end{proof}

For the Hahn field in \eqref{eq:hahn}, $\OO_F$ consists of series of
nonnegative valuation, and standard part extracts the coefficient of $t^0$.
There is no field homomorphism $F\to\R$ extending this map when $F$ has a
nonzero infinitesimal: $\st(\varepsilon)=0$ would conflict with
$\varepsilon\varepsilon^{-1}=1$. Infinite curvature therefore has no finite
real standard part simply because it has acquired a surreal name.

An observational rule based on standard part is a possible modeling choice,
not a consequence of the field axioms. In particular, it discards infinitesimal
corrections by design. A theory that instead measures them needs to specify what
apparatus, outcomes, and probability rule do so.

\subsection{Transfer: what one may and may not infer}
Real-closed-field algebra has first-order transfer in the language of ordered
fields. Richer elementary embeddings preserving restricted analytic and
exponential structure also exist in the settings studied by Ehrlich and Kaplan
\cite{EK}. It would be incorrect to say that surreals admit no transfer.

What does not follow from a bare field embedding is a transfer principle for all
real functions, all subsets, Lebesgue measure, Hilbert spaces, or theorems about
solutions of nonlinear partial differential equations. The language, the
expanded structure, and the admissible quantifiers matter. A nonstandard
superstructure carrying internal sets and transferred integration supplies
additional data that the ordered field alone does not specify. Thus ``the
formula still makes sense'' is a much weaker statement than ``the existence,
uniqueness, and stability theorem still holds.''

\section{Topology and calculus are part of the physical proposal}\label{sec:calculus}

\subsection{Why the full fine topology is not ordinary spacetime topology}
Use the positive surreal-valued modulus to define fine balls in $\Sc$:
$B(a,r)=\{z:|z-a|<r\}$, with every positive surreal $r$ allowed. The following
simple argument explains a central warning in the repository
\cite{RepoAnalysis}.

\begin{proposition}[Set-sized collapse in the full fine topology]\label{prop:discrete}
Every set $A\subseteq\Sc$ is closed and discrete in the fine topology. Every
convergent set-indexed net is eventually constant. A continuous map from the
ordinary real interval $[0,1]$ to this fine space is constant.
\end{proposition}
\begin{proof}
For a fixed $b$, the nonzero distances $|a-b|$, with $a\in A$, form a set of
positive surreal numbers. If it is nonempty, the Conway cut
$\{0\mid |a-b|:a\in A,\ a\ne b\}$ gives a positive number smaller than every
one of these distances. The resulting ball isolates $b$ from all the other
points of $A$; if the distance set is empty, any positive radius suffices.
For $b\notin A$ this also proves that the complement of $A$ is open.
The range of a set-indexed net, together with its proposed limit, is a set, so a
ball isolating the limit forces eventual equality. Finally the image of
$[0,1]$ is a set and hence discrete; a connected interval has no nonconstant
continuous map into a discrete space.
\end{proof}

This is not a theorem that all useful surreal analysis is impossible. It is a
theorem that a particular full-class topology is unsuitable for importing
ordinary real-parameter continuous motion without modification. It also applies
to the inclusion of the real line: its usual topology is not the subspace
topology inherited from this full fine space.

A fixed Hahn field has an intrinsic order or valuation topology. For a rank-one
value group, $t^n\to0$ in its valuation topology when $\val(t)>0$. This does not
contradict \cref{prop:discrete}: in the full surreal field there are additional
positive radii smaller than every $t^n$. The topology on a fixed workspace and
the topology induced from the entire class need not coincide.

A standard-part topology on finite elements, defined by inverse images of
ordinary open sets under $\st$, goes in the opposite direction. It cannot
topologically separate two points with the same standard part. Its natural
separated quotient is ordinary space. A coefficientwise topology is yet another
choice. None should be silently substituted for another in a proof about causal
curves or limits.

\subsection{Three derivatives that must not be conflated}
Consider a field with coefficients $a_\gamma(x)$ depending on ordinary spacetime
coordinates $x=(x^0,\ldots,x^3)$. A natural derivative is
\begin{equation}\label{eq:coeffder}
\partial_\mu\left(\sum_\gamma a_\gamma(x)t^\gamma\right)
 =\sum_\gamma (\partial_\mu a_\gamma)(x)t^\gamma,
\qquad \partial_\mu t^\gamma=0.
\end{equation}
This is spacetime differentiation with the asymptotic parameter held fixed.
A derivative with respect to that parameter, such as
$\dd(\varepsilon^q)/\dd\varepsilon=q\varepsilon^{q-1}$, is a different
operation. A derivation on the scalar field $\No$ is different again.

The distinguished Berarducci--Mantova derivation $D$ kills $\R$, is compatible
with the surreal real exponential, and can be normalized by $D\omega=1$.
Consequently
\begin{equation}\label{eq:BMscale}
D(\omega^{-1})=-\omega^{-2}=-\varepsilon^2,
\end{equation}
not $1$ and not $0$ \cite{BM}. There is no contradiction: the three operations
answer different questions. In particular, $D$ does not act as an ordinary time
derivative on the real number $t$, because $D(t)=0$ for every real scalar $t$.
Nor does one scalar derivation automatically supply four independent commuting
coordinate derivatives and a covariant connection.

\subsection{There already are substantial function and integration theories}
Costin and Ehrlich construct extensions of a specified class of resurgent real
functions to surreal arguments, together with compatible antidifferentiation
and integration. Their class includes important Borel-summable examples and
solutions of nonresonant meromorphic ODE systems; changes of variables allow
finite singular locations to be considered. They also prove obstructions to
substantially more unrestricted constructive extensions \cite{CE}.

This is a positive starting point for physics. A reduced gravitational ODE
might fall within such a class, and then a surreal function extension would be
more informative than simply assigning a formal value to its first divergent
term. Membership, the relevant sector or branch, and the chosen integration
operator would have to be checked. These results do not amount to a general
Einstein PDE existence theorem or a complete theory of oscillatory quantum
fields on a surreal spacetime. They concern derivations on functions, not the
same object as the scalar derivation in \eqref{eq:BMscale}.

\subsection{Common domains and support conditions are essential}
The expression
\[
A(x)=\sum_\gamma a_\gamma(x)t^\gamma
\]
will be used below only when the coefficient functions live on a specified
common ordinary domain and the support rules justify the operations being
performed. Having an analytic germ for each coefficient is weaker than having a
common neighborhood. For example,
\[
\sum_{n\ge1}\frac{t^n}{1-nz}
\]
has individually analytic coefficients at zero but no ordinary disk on which
all of them are holomorphic. This is a useful warning from the repository's
ring comparison \cite{RepoDocs}.

Likewise, a Taylor expansion does not uniquely determine an arbitrary smooth
function. The ordinary function $\ee^{-1/x^2}$ for $x>0$, continued by zero for
$x\le0$, has zero Taylor series at zero and is not identically zero. A rule that
extends all smooth functions by formal Taylor coefficients necessarily loses
such information. A physical proposal must choose a function class or a richer
extension rule rather than assuming that every classical field automatically
extends canonically.

\section{What a black-hole singularity is}\label{sec:singularity}

\subsection{Coordinate divergence, curvature divergence, and incompleteness}
A divergent coordinate component need not be singular geometry. A curvature
invariant can distinguish many genuine divergences from this coordinate
artifact, but absence of a particular divergent invariant is not a general
regularity theorem. Most fundamentally, an inextendible causal geodesic may
have finite affine length, or finite proper time in the timelike case. Whether
the spacetime has an extension then depends on the allowed regularity and field
equations \cite{Witten,Dafermos,DL}.

These distinctions survive any serious change of scalars. A metric needs to be
nondegenerate to define its inverse. Curvature needs derivatives in a specified
sense. A geodesic needs an affine parameter and a solution theory. An extension
needs a larger space, compatible charts or their replacement, and agreement
with the original geometry. Adding $\omega$ to the scalar vocabulary supplies
none of these structures by itself.

\subsection{What the singularity theorems actually constrain}
A standard form of Penrose's theorem assumes global hyperbolicity with a
noncompact Cauchy surface, an appropriate compact trapped surface, and the null
convergence condition. It yields future null geodesic incompleteness
\cite{Penrose}. The conclusion is not that every possible singularity is a
point with infinite density, or that every curvature scalar must diverge.

A new theory may evade a hypothesis: it might modify the effective energy
condition, causal structure, topology, differentiability, or spacetime concept.
However, merely adopting an ordered extension of $\R$ does not identify which
hypothesis changes or prove a replacement theorem. Conversely, one cannot
blindly transfer the entire classical proof to a fine-topological surreal
manifold, because global compactness and causal-topological arguments have not
been preserved.

\subsection{Focusing illustrates the issue without settling the whole theorem}
For a twist-free null congruence in four dimensions, the Raychaudhuri equation
has the form
\begin{equation}\label{eq:ray}
\frac{\dd\theta}{\dd\lambda}
 =-\frac12\theta^2-\sigma_{ab}\sigma^{ab}-R_{ab}k^ak^b.
\end{equation}
Here the shear contraction is on the positive screen space. Under null
convergence, the right side is at most $-\theta^2/2$. The equality model is
\begin{equation}\label{eq:raymodel}
\theta(\lambda)=\frac{\theta_0}{1+\tfrac12\theta_0\lambda},
\qquad \theta_0<0,
\end{equation}
with a pole at $\lambda_*=-2/\theta_0$ \cite{Raychaudhuri,Witten}.

Evaluating \eqref{eq:raymodel} at $\lambda_* -\varepsilon$ replaces a large
negative real expansion by an infinite negative surreal expansion. At
$\lambda_*$ its denominator is still zero. This illustrates why infinite
numbers alone do not prevent focusing. A caustic of a congruence is not itself
necessarily a spacetime singularity; the global argument is what turns
focusing, under further assumptions, into incompleteness.

\subsection{Changing variables does not automatically change the physics}
For the elementary equation $y'=y^2$, the solution $y=(t_*-t)^{-1}$ blows up.
The new variable $u=1/y$ obeys $u'=-1$ and crosses zero smoothly. This is a
regular equation for $u$, but the inverse map fails at $u=0$. Whether $u$ or $y$
is the physical observable is additional information. Likewise,
$s=-\log(\tau_* -\tau)$ sends a finite-proper-time endpoint to infinite
coordinate time without changing the elapsed proper time. A compactification,
a blow-up, a reciprocal field, and a geodesic completion are different
operations.

\section{Schwarzschild at an infinitesimal radius}\label{sec:schwarzschild}

\subsection{Units and the event horizon}
Set $G=c=1$ for the calculations. The mass parameter $m$ has dimensions of
length; in physical units $m=GM/c^2$. Proper time $\tau$ is expressed in length
units, so that $\tau=c\,\tau_{\rm physical}$. A dimensionless ratio, rather
than an unqualified dimensionful number, will be declared infinitesimal.

The Schwarzschild metric is
\begin{equation}\label{eq:schwarz}
\dd s^2=-f(r)\dd t^2+f(r)^{-1}\dd r^2+r^2\dd\Omega^2,
\qquad f(r)=1-\frac{2m}{r},
\end{equation}
where $\dd\Omega^2=\dd\theta^2+\sin^2\theta\,\dd\phi^2$. Introduce the
advanced coordinate $v=t+r_*$, with $\dd r_*/\dd r=f^{-1}$. Then
\begin{equation}\label{eq:ef}
\dd s^2=-f(r)\dd v^2+2\dd v\,\dd r+r^2\dd\Omega^2.
\end{equation}
The $v,r$ block has determinant $-1$ even at $r=2m$. Away from the ordinary
angular-coordinate poles, the full determinant is $-r^4\sin^2\theta$.
Thus the divergence of $g_{rr}$ at the horizon in \eqref{eq:schwarz} is already
removed by an ordinary real coordinate change. No generalized number is needed.

\subsection{The invariant calculation}
For a metric of the form \eqref{eq:schwarz} with general $f$, direct curvature
calculation gives
\begin{equation}\label{eq:Kgeneral}
\Kret[f]:=R_{abcd}R^{abcd}
 =(f'')^2+4\left(\frac{f'}r\right)^2
       +4\left(\frac{1-f}{r^2}\right)^2.
\end{equation}
A derivation is supplied in \cref{app:curvature}. Substituting $f=1-2m/r$ gives
\begin{equation}\label{eq:Kschwarz}
\Kret=\frac{48m^2}{r^6},\qquad
\Kret(2m)=\frac{3}{4m^4}.
\end{equation}
Let $x=r/m$. Then the dimensionless invariant is
\begin{equation}\label{eq:Keps}
m^4\Kret=48x^{-6};\qquad x=\varepsilon>0
\quad\Longrightarrow\quad m^4\Kret=48\varepsilon^{-6}.
\end{equation}
This is an exact and useful non-Archimedean description. It records both the
power and coefficient of the divergence. It is larger than every positive real
number and has no finite standard part. At $r=0$, the defining expression still
has a pole.

For $r>0$, Schwarzschild is vacuum: its Ricci tensor and Ricci scalar vanish.
The divergent invariant is built from the full Riemann tensor, equivalently the
Weyl tensor in this vacuum region. It is consequently misleading to explain
this example as simply an infinite ordinary matter density that can be replaced
by a surreal density. The explicit Ricci check appears in the companion
verification program.

\subsection{A precise pole-persistence statement}
\begin{proposition}[Poles survive scalar extension and ramification]\label{prop:pole}
Let $p\ge1$ be an integer, let $u(r)\in\R[[r]]$ have nonzero constant term
$c$, and put $f(r)=r^{-p}u(r)$. Evaluate the series at a positive infinitesimal
$\varepsilon$ in a Hahn field, using strong summation. Then
\[
\val f(\varepsilon)=-p\,\val\varepsilon,
\qquad f(\varepsilon)=c\varepsilon^{-p}(1+\delta),
\quad \delta\in\mm_F.
\]
In particular $f(\varepsilon)$ is infinite in magnitude. Under a formal change
$r=s^q w(s)$, with $q\ge1$ an integer and $w(0)\ne0$, the pole order becomes
$pq$ and does not disappear.
\end{proposition}
\begin{proof}
Strong summation of $u(\varepsilon)$ is valid by the positive-valuation power
series rule. Its constant term is $c$, so $u(\varepsilon)=c(1+\delta)$ with
$\val\delta>0$, unless $\delta=0$. Multiplication by
$\varepsilon^{-p}$ gives the asserted leading term and valuation. For the
change of variable,
\[
f(s^q w(s))=s^{-pq}\,w(s)^{-p}u(s^q w(s)),
\]
and the factor following $s^{-pq}$ is a unit with constant term
$w(0)^{-p}c$. Its nonzero constant term cannot cancel the pole.
\end{proof}

The corresponding real-analytic statement has a direct geometric consequence.
If a dimensionless scalar curvature invariant tends to infinity along a curve,
there is no extension along that curve to an ordinary nondegenerate $C^2$
metric at an endpoint where that invariant is continuous and finite: continuity
would give a finite limit. This is a restricted and useful obstruction, not a
claim that all conceivable generalized geometries have been ruled out.

\subsection{An infinite number is not an added regular event}
Assigning a nonzero infinitesimal radius describes another point in a punctured
extension of the coordinate domain. It does not add the missing endpoint.
There is also no ``last positive radius'' before zero: $\varepsilon/2$ is
smaller. A proposed rule that replaces zero by one chosen infinitesimal has
introduced a regulator. It must explain the choice, its physical units, the
resulting equations, and invariance of observations under changing that
choice.

Allowing a metric to have infinite coefficients might be part of a broader
theory. In that case the theory must define regularity for those coefficients
and the behavior of clocks, rods, and tidal observables there. It cannot infer
regularity solely from the fact that a coefficient belongs to $\No$.

\section{Proper-time infall and the physical tidal field}\label{sec:infall}

\subsection{An exact radial geodesic}
For a radial timelike Schwarzschild geodesic with energy per unit rest mass
$E=1$, corresponding to rest at infinity, the first integral is
\begin{equation}\label{eq:infall}
\left(\frac{\dd r}{\dd\tau}\right)^2=E^2-f(r)=\frac{2m}{r}.
\end{equation}
Choose the inward sign and integrate:
\[
r^{1/2}\dd r=-\sqrt{2m}\,\dd\tau,
\qquad
\frac23r^{3/2}=\sqrt{2m}(\tau_*-\tau).
\]
Writing $u=\tau_*-\tau>0$ gives the exact relation
\begin{equation}\label{eq:properr}
r^3=\frac92m u^2,
\qquad
\Kret=\frac{64}{27u^4}.
\end{equation}
Thus the singular behavior occurs at a finite proper time. The fractional power
$r\propto u^{2/3}$ is particularly well suited to a ramified Hahn or Puiseux
workspace. This is an improvement in representation, not in the boundedness of
\eqref{eq:properr}.

\subsection{Geodesic deviation resolves the scale hierarchy}
In a parallel-propagated radial freely falling orthonormal frame, the radial
and transverse tidal equations are
\begin{equation}\label{eq:jacobi}
\frac{\dd^2\xi_r}{\dd\tau^2}=\frac{2m}{r^3}\xi_r,
\qquad
\frac{\dd^2\xi_\perp}{\dd\tau^2}=-\frac{m}{r^3}\xi_\perp.
\end{equation}
Substituting \eqref{eq:properr} changes their coefficients to
$4/(9u^2)$ and $-2/(9u^2)$. Because
$\dd^2/\dd\tau^2=\dd^2/\dd u^2$, the ansatz $\xi=u^a$ produces the
indicial equations
\[
a(a-1)=\frac49\quad\hbox{or}\quad a(a-1)=-\frac29.
\]
Their solutions are
\begin{align}
\xi_r&=A u^{4/3}+B u^{-1/3},\label{eq:radmodes}\\
\xi_\perp&=C u^{2/3}+D u^{1/3}.\label{eq:transmodes}
\end{align}
For a generic diagonal family of three independent Jacobi separations, the
radial direction stretches as $u^{-1/3}$ while the two transverse directions
shrink as $u^{1/3}$. The associated volume scales as $u^{1/3}$ and tends to
zero. Vanishing leading coefficients only changes the indicated dominant
powers, not the divergence of the curvature itself.

The surreal description retains all these scales simultaneously. For
$u/m=\varepsilon$, it records an infinite tidal coefficient, an infinite radial
stretching mode, and infinitesimal transverse separations. Calling these
legitimate numbers is mathematically coherent. Calling them bounded physical
forces or a nonsingular passage is not justified.

\subsection{Why a formal bounce is not a geodesic extension}
One can write $r\propto |u|^{2/3}$ on both sides of $u=0$. This operation does
not produce a smooth Schwarzschild continuation through the singularity. The
first derivative of $r$ is unbounded there, the curvature still behaves as
$u^{-4}$, and the original metric and geodesic equation are undefined at the
joining event. Matching two formulas across an omitted point does not supply
the missing geometry or a unique law for outgoing matter.

There is no assertion here that a quantum-gravitational bounce is impossible.
Such a bounce would need a different dynamical or geometric account of the
transition. Surreal methods might describe its asymptotic incoming and outgoing
regions, but the transition law is precisely the physical information that the
classical formulas do not contain.

\section{A useful positive theorem: formal Einstein reduction}\label{sec:einstein}

\subsection{A conservative geometric framework}
A practical first model does not replace the set of spacetime events. Let $U$
be an ordinary coordinate domain in a smooth four-dimensional manifold. Attach
to it the ring
\[
\mathcal A(U)=C^\infty(U,\R)[[h]],
\]
where $h$ is a formal parameter fixed by the coordinate derivatives. Fields are
formal families on an ordinary spacetime, not arbitrary functions on a
fine-topological surreal manifold. Formal power series can later be evaluated
at a positive infinitesimal in a suitable Hahn workspace, coefficient by
coefficient.

The notation deliberately separates two questions: whether a formal equation
is well defined, and whether its formal solution represents an actual family
of real solutions with controlled remainders. The first is algebraic. The
second is an analytic problem and generally harder.

\begin{theorem}[Reduction of the Einstein tensor]\label{thm:einstein}
Let
\[
g(h)=g_0+h g_1+h^2g_2+\cdots
\]
be a symmetric $4\times4$ matrix with entries in $\mathcal A(U)$, and assume
$g_0$ is an ordinary smooth nondegenerate Lorentzian metric on $U$. Define the
connection and curvature by the usual coordinate formulas, differentiating
coefficients and holding $h$ fixed. Then the inverse metric exists in
$M_4(\mathcal A(U))$, and reduction $\pi_0$ modulo $h$ satisfies
\begin{align}
\pi_0(g^{-1})&=g_0^{-1},\label{eq:ginvreduce}\\
\pi_0\bigl(\Ric[g(h)]\bigr)&=\Ric[g_0],\label{eq:ricreduce}\\
\pi_0\bigl(G[g(h)]\bigr)&=G[g_0].\label{eq:einreduce}
\end{align}
If $\Lambda(h)$ and $T(h)$ have the same coefficientwise interpretation and
\[
G[g(h)]+\Lambda(h)g(h)=8\pi T(h),
\]
then
\[
G[g_0]+\Lambda_0g_0=8\pi T_0.
\]
\end{theorem}
\begin{proof}
Write $g=g_0(I+H)$ with $H\in hM_4(\mathcal A(U))$. Then
\[
g^{-1}=\left(\sum_{n\ge0}(-H)^n\right)g_0^{-1}.
\]
At any specified power of $h$, only finitely many terms contribute. This is a
formal inverse, so no topological convergence is being used. It proves
\eqref{eq:ginvreduce}.

The map $\pi_0$ preserves finite sums and products and commutes with
$\partial_\mu$, since the latter differentiates only coefficient functions.
Apply it to
\[
\Gamma^a{}_{bc}
=\frac12 g^{ad}\bigl(\partial_b g_{dc}+\partial_c g_{db}
                         -\partial_d g_{bc}\bigr)
\]
and then to
\[
R^a{}_{bcd}=\partial_c\Gamma^a{}_{db}-\partial_d\Gamma^a{}_{cb}
 +\Gamma^a{}_{ce}\Gamma^e{}_{db}-\Gamma^a{}_{de}\Gamma^e{}_{cb}.
\]
Reduction commutes with the contractions defining Ricci curvature and scalar
curvature. It therefore commutes with
$G_{ab}=R_{ab}-\tfrac12\Scal(g)g_{ab}$. Applying reduction to the formal
Einstein equation proves its last assertion.
\end{proof}

The same argument works for common-domain Hahn coefficient families when their
supports justify multiplication and the inverse expansion. It is essential to
state those hypotheses rather than infer them from pointwise values. The
formal theorem is self-contained; it is not a claim that the repository has
already formalized Einstein geometry.

\subsection{What the theorem says about new physics}
A regular infinitesimal deformation of a nondegenerate classical metric obeys
the classical Einstein equation at leading order. Thus a model cannot both keep
this regular common-domain reduction and obtain an arbitrary different
macroscopic leading equation without changing another assumption, such as the
source, the action, or the scaling regime.

The theorem does \emph{not} exclude singular perturbations. A leading metric may
become degenerate at a boundary, inverse coefficients may acquire negative
valuation, or derivatives may bring a seemingly small correction to leading
order. A rescaled inner problem can then differ substantially from the outer
problem. Those are exactly the situations in which a controlled asymptotic
language is useful. They are not counterexamples to the theorem, whose
nondegeneracy and common-domain assumptions have ceased to apply.

\subsection{Why a pointwise standard part is insufficient}
For a positive ordinary real parameter $\delta$, consider
\begin{equation}\label{eq:boundarylayer}
b_\delta(x)=\frac{\delta^2}{1+(x/\delta)^2}
           =\frac{\delta^4}{\delta^2+x^2}.
\end{equation}
For every real $x$, $0\le b_\delta(x)\le\delta^2$, so the functions converge
uniformly to zero. However,
\begin{equation}\label{eq:boundarylayerder}
b_\delta''(0)=-2.
\end{equation}
The limit of the functions has second derivative zero. Uniform smallness of a
metric perturbation therefore does not by itself imply small curvature. This
particular family even has first derivatives uniformly of order $\delta$;
second-derivative control is the missing ingredient.

The same rational formula can be evaluated at an infinitesimal $\delta$.
Every pointwise value has zero standard part, but a coefficientwise or
appropriately extended second derivative at zero still yields $-2$.
There is no common formal expansion in nonnegative powers of $\delta$ with
smooth coefficients valid across $x=0$: away from zero, expansion in
$\delta^2/x^2$ produces coefficients singular at zero. Thus the example does
not contradict \cref{thm:einstein}. It demonstrates exactly why its hypotheses
cannot be omitted.

\subsection{Actions and conservation laws, with an honest scope}
On a fixed ordinary compact integration domain, a common-support family of
smooth densities can be integrated coefficientwise:
\[
\int_U\left(\sum_\gamma L_\gamma(x)t^\gamma\right)\dd^4x
 :=\sum_\gamma\left(\int_U L_\gamma(x)\dd^4x\right)t^\gamma.
\]
The support of the result is a subset of the original well-ordered support.
Ordinary integration by parts applies separately to every coefficient, under
its usual boundary assumptions. Consequently a formal local action principle,
Euler--Lagrange equations, and coefficientwise conservation identities can be
developed without pretending that a fine-topological Riemann integral over the
whole surreal class has been constructed.

For the Einstein--Hilbert density, the inverse metric and the branch of
$\sqrt{-\det g}$ exist formally around a nondegenerate Lorentzian $g_0$.
This makes a formal gravitational action feasible. It does not define a
nonperturbative path-integral measure, select boundary conditions at a
singularity, or prove convergence of the expansion.

\section{A regular-core model: what changes, and what does not}\label{sec:core}

\subsection{A real modification is already enough to remove the central pole}
To separate a change of dynamics from a change of scalars, consider the
static metric function used in Hayward's regular-core model
\cite{Hayward}:
\begin{equation}\label{eq:hayward}
f_\ell(r)=1-\frac{2mr^2}{r^3+2m\ell^2},\qquad m>0,\quad\ell>0.
\end{equation}
Here $\ell$ is an ordinary real length scale for now. For large $r$ the metric
approaches Schwarzschild, whereas at the center
\begin{equation}\label{eq:haywardcenter}
f_\ell(r)=1-\frac{r^2}{\ell^2}
              +\frac{r^5}{2m\ell^4}+O(r^8).
\end{equation}
Inserting this expansion into \eqref{eq:Kgeneral} gives
\begin{equation}\label{eq:haywardK}
\Kret(0)=\frac{24}{\ell^4}.
\end{equation}
The local de Sitter-type core and finite limiting curvature are consequences of
the \emph{new metric function}, not of any non-Archimedean arithmetic. The
vanishing spherical-coordinate determinant at a center must be assessed in
Cartesian coordinates; the expansion above yields a nondegenerate $C^2$ local
metric there. No claim about all higher derivatives follows from the curvature
limit alone.

The effective mass function is
\[
M_\ell(r)=\frac{mr^3}{r^3+2m\ell^2},
\qquad f_\ell=1-\frac{2M_\ell(r)}r.
\]
Using $G^t{}_t=-2M_\ell'(r)/r^2$ and $G_{ab}=8\pi T_{ab}$ yields the effective
energy density
\begin{equation}\label{eq:haywardrho}
\rho(r)=\frac{M_\ell'(r)}{4\pi r^2}
       =\frac{3m^2\ell^2}{2\pi(r^3+2m\ell^2)^2},
\qquad \rho(0)=\frac{3}{8\pi\ell^2}.
\end{equation}
Also $p_r=-\rho$; at the center the stress is of cosmological-constant form.
These identities exhibit the changed source explicitly. Writing
$T_{ab}=G_{ab}/8\pi$ is a consistent effective description, but it does not by
itself provide a microphysical matter Lagrangian or establish stability.

A finite-curvature core alone is not a proof of global geodesic completeness,
unique evolution, or physically stable evaporation. The regularity class,
inner-horizon behavior, global causal assumptions, and stress-energy
interpretation must all be examined. Classifications of geodesically complete
black-hole scenarios make precisely this larger geometric issue explicit
\cite{CR}.

\subsection{Making the core radius infinitesimal changes the conclusion}
Now set $\ell=m\varepsilon$ with a positive infinitesimal $\varepsilon$.
For every fixed positive ordinary regulator the curvature in
\eqref{eq:haywardK} is finite. As a non-Archimedean value, however,
\[
m^4\Kret(0)=24\varepsilon^{-4}
\]
is infinite. The scalar field has not supplied a finite observational curvature
bound. There are two different statements:
\emph{regular for each fixed real regulator} and \emph{uniformly finite in the
regulator's infinitesimal limit}. The first does not imply the second.

Keeping a finite real $\ell$, possibly motivated by a physical ultraviolet
scale, can give finite central invariants. Choosing the scale and the modified
source remains a physical hypothesis. Surreal arithmetic does not determine
that $\ell$ is the Planck length, or that a particular regular-core ansatz is
the correct quantum-gravitational dynamics.

\subsection{Hahn summability detects the correct boundary layer}
The model nevertheless supplies a particularly clear positive application of
Hahn methods. Write
\begin{equation}\label{eq:dimensionlesscore}
x=\frac rm,\qquad \varepsilon=\frac\ell m,
\qquad f(x,\varepsilon)=1-\frac{2x^2}{x^3+2\varepsilon^2}.
\end{equation}
Treat $\varepsilon$ as a Hahn monomial of positive valuation and take
$x=\varepsilon^\alpha$ with $\alpha>0$ real. The requisite real exponents can
be included in a divisible set-sized value group.

\begin{proposition}[Outer and inner support criteria]\label{prop:matching}
For \eqref{eq:dimensionlesscore}, the outer geometric expansion
\begin{equation}\label{eq:outer}
f=1-\frac2x\sum_{k\ge0}
                  \left(-\frac{2\varepsilon^2}{x^3}\right)^k
\end{equation}
is strongly Hahn summable when $\alpha<2/3$, whereas the inner expansion
\begin{equation}\label{eq:inner}
f=1-\frac{x^2}{\varepsilon^2}\sum_{k\ge0}
                  \left(-\frac{x^3}{2\varepsilon^2}\right)^k
\end{equation}
is strongly Hahn summable when $\alpha>2/3$.
At $\alpha=2/3$, neither displayed geometric family is strongly summable.
The exact transition-scale expression, with $x=\varepsilon^{2/3}X$, is
\begin{equation}\label{eq:transition}
f=1-2\varepsilon^{-2/3}\frac{X^2}{X^3+2}.
\end{equation}
For positive finite $X$ with positive real standard part, its denominator is a
unit and the expression is well defined.
\end{proposition}
\begin{proof}
The ratio in \eqref{eq:outer} has valuation
$(2-3\alpha)\val(\varepsilon)$. It is positive exactly for $\alpha<2/3$,
so its powers have increasing well-ordered support and the geometric identity
applies. If $\alpha>2/3$, their exponents strictly decrease and the family is
not Hahn summable. The ratio in \eqref{eq:inner} has valuation
$(3\alpha-2)\val(\varepsilon)$, giving the complementary criterion.
At equality all geometric terms contribute at the same exponent; infinitely
many nonzero contributions to that exponent violate strong summability.
Equation \eqref{eq:transition} follows by direct substitution in the original
rational function, without summing either invalid geometric family.
\end{proof}

For every finite $N$, the remainder identity underlying the proof is exact:
\begin{equation}\label{eq:remainder}
\frac1{1+A}=\sum_{k=0}^{N-1}(-A)^k+\frac{(-A)^N}{1+A}.
\end{equation}
If $\val A>0$, the remainder has valuation $N\val A$. This provides a
certificate of asymptotic order and a natural target for the repository's
existing Neumann-series and rational-jet infrastructure.

The crossover radius is $r^3\sim2m\ell^2$. A formal small correction in the
outer region becomes leading in the inner region. The calculation illustrates
an actual gain: support conditions detect when an approximation has been used
outside its asymptotic regime. At the transition the failure is a failure of
that expansion, not a singularity of the rational model. No new physical
solution has been invented by the series algebra.

\section{Effective-field-theory scales and more general singular interiors}\label{sec:eft}

\subsection{The Planck scale is a positive real scale, not an infinitesimal}
Low-energy quantum gravity can be treated as an effective field theory; its
controlled use depends on suitable separation from the ultraviolet scale, not
on an assumption that the classical equations hold at arbitrarily high
curvature \cite{Donoghue,DonoghueReview}. A simple curvature diagnostic uses
\[
\ell_P=\sqrt{\frac{\hbar G}{c^3}},
\qquad \ell_P^4\Kret.
\]
On the Schwarzschild solution, the condition $\ell_P^4\Kret\sim1$ gives
\begin{equation}\label{eq:rq}
r_Q\sim48^{1/6}(m\ell_P^2)^{1/3}.
\end{equation}
This is dimensional power counting, not a derived universal threshold or an
exact prediction of a quantum-gravity theory. Other state-dependent scales and
high frequencies can also matter. It is nevertheless enough to show the
logical problem with extrapolating unchanged classical equations all the way
to an infinitesimal radius while keeping $m$ and $\ell_P$ ordinary and fixed:
a positive real high-curvature scale is encountered first.

For $m\gg\ell_P$, horizon curvature can be small even though the interior
ultimately becomes a high-curvature region. A black-hole horizon, a
Planck-curvature regime, and a classical singular boundary are distinct.

\subsection{Two-parameter power counting becomes exact valuation arithmetic}
One can organize a family in which both ratios vary. Let
\[
\frac rm=\varepsilon^\alpha,
\qquad \frac{\ell_P}{m}=\varepsilon^\beta,
\qquad \alpha,\beta>0.
\]
Then
\begin{equation}\label{eq:EFTvaluation}
\ell_P^4\Kret=48\varepsilon^{4\beta-6\alpha}.
\end{equation}
Thus the curvature parameter is infinitesimal if $2\beta>3\alpha$, of order
one if $2\beta=3\alpha$, and infinite if $2\beta<3\alpha$. This separates
three asymptotic regimes without conflating the small-body, large-mass, and
near-singularity limits. The ratio assignments describe a family of models;
they do not assert that a fixed measured Planck length is literally an
infinitesimal.

Valuation arithmetic is especially attractive when several corrections compete.
For example a candidate term proportional to
$\varepsilon^a x^{-b}$ changes its order according to $a-b\alpha$. A
collection of such inequalities can locate dominant balances and propose
rescalings before a costly symbolic expansion is attempted. The physical
coefficients and allowed operators still have to come from the effective
action or another specified dynamical model.

\subsection{Kasner behavior and the value of fractional powers}
A second exact test is the vacuum Kasner metric, with dimensionful reference
scales absorbed into the coordinates:
\begin{equation}\label{eq:kasner}
\dd s^2=-\dd\tau^2+\sum_{j=1}^3\tau^{2p_j}\dd x_j^2,
\qquad \sum_j p_j=1,\quad \sum_jp_j^2=1.
\end{equation}
Its curvature is
\begin{equation}\label{eq:kasnerK}
\Kret=4\tau^{-4}\left[
 \sum_j p_j^2(p_j-1)^2+\sum_{j<k}p_j^2p_k^2\right].
\end{equation}
These expressions follow directly from a diagonal orthonormal frame; their
constraints and special cases are checked in the companion program. For
$(p_1,p_2,p_3)=(-1/3,2/3,2/3)$ the coefficient is $64/27$; for
$(1,0,0)$ it vanishes, corresponding to a flat coordinate example rather than
a genuine curvature divergence.

Hahn monomials efficiently record anisotropic scale factors and corrections to
a Kasner regime. Rigorous singular asymptotics already exist in special
Einstein--matter settings: Andersson and Rendall's quiescent Einstein--scalar
and stiff-fluid analysis is a concrete benchmark \cite{AR}. A useful surreal
project would reproduce a stated solution class and its error estimates, not
replace those estimates by an unsupported appeal to infinitesimal notation.
Generic oscillatory sequences of changing Kasner epochs present additional
issues; a single well-ordered expansion must not be assumed to encode every
possible interior evolution.

\subsection{Rotating and charged interiors are not Schwarzschild copies}
Inner-horizon instability and mass inflation in charged or rotating models
introduce blueshift, late-time tails, and regularity questions different from
the spacelike Schwarzschild singularity. Dafermos's charged-black-hole analysis
and the conditional interior result of Dafermos and Luk distinguish geometric
extension from stronger regularity and uniqueness questions \cite{Dafermos,DL}.
In particular, a $C^0$ extension does not by itself provide finite curvature or
a classical $C^2$ solution beyond the boundary.

Exponential growth factors together with inverse powers or logarithms of an
advanced-time parameter are natural candidates for transseries descriptions.
Surreal exponentiation and logarithms can retain separated scales that a
truncated power series misses. But an expression schematically involving
$\ee^{\kappa v}v^{-p}$ is not a universal theorem about all Kerr interiors;
the model, tail assumptions, gauge, and quantities being estimated must be
specified. Computing a large generalized scalar is not a substitute for the
nonlinear stability argument.

\section{Surcomplex amplitudes: what finite-dimensional quantum algebra permits}\label{sec:quantum}

\subsection{The algebraic part works well}
Let $F\supseteq\R$ be a real-closed ordered field and let $K=F[i]$. Define
\[
\overline{a+ib}=a-ib,\qquad |a+ib|=\sqrt{a^2+b^2},
\qquad \langle\psi,\phi\rangle=\sum_{j=1}^n\overline{\psi_j}\phi_j.
\]
This gives a positive-definite Hermitian form on the finite-dimensional space
$K^n$. Orthogonality, unitary matrices, tensor products of finitely many systems,
and projective-state algebra all make sense.

The same finite algebra supports Lorentzian quadratic forms, their linear
isometry groups, Clifford algebras, and finite spinor constructions over
appropriate extensions. What has been preserved here is an algebraic
structure. Compactness of a group, spectral measures, the topology of a state
space, and continuum dynamics are separate issues.

\begin{theorem}[Finite Hermitian spectral theorem]\label{thm:spectral}
Every Hermitian matrix $A=A^\dagger\in M_n(K)$ has an orthonormal eigenbasis,
and all its eigenvalues belong to $F$.
\end{theorem}
\begin{proof}
The characteristic polynomial splits over the algebraically closed field $K$,
so $A$ has an eigenvector $v\ne0$ with eigenvalue $\lambda$. Since
$v^\dagger v>0$,
\[
\lambda v^\dagger v=v^\dagger Av
 =\overline{v^\dagger Av}=\overline\lambda\,v^\dagger v,
\]
and hence $\lambda=\overline\lambda\in F$. Normalize $v$ using the positive
square root of $v^\dagger v$. The orthogonal complement of $v$ is invariant
under $A$, because
$\langle v,Aw\rangle=\langle Av,w\rangle=\lambda\langle v,w\rangle$.
Induction on $n$ completes the proof.
\end{proof}

This proof requires no completeness of $F$. It is a finite algebraic theorem,
not an infinite-dimensional spectral theorem for unbounded operators.

\subsection{A precise observational-shadow result}
\begin{theorem}[Finite quantum experiments admit ordinary shadows]\label{thm:quantumshadow}
Let $\psi\in K^n$ satisfy $\psi^\dagger\psi=1$. Its components are finite,
and $\psi_0=\st\psi\in\C^n$ satisfies $\psi_0^\dagger\psi_0=1$.
Every unitary matrix $U\in M_n(K)$ has finite entries, and $U_0=\st U$ is
unitary. Every orthogonal projection $P=P^\dagger=P^2$ has finite entries, and
$P_0=\st P$ is an ordinary orthogonal projection. Moreover,
\begin{equation}\label{eq:bornshadow}
\st\bigl(\psi^\dagger P\psi\bigr)=\psi_0^\dagger P_0\psi_0,
\qquad \st(U\psi)=U_0\psi_0.
\end{equation}
Consequently the real standard parts of joint probabilities in a fixed finite
sequence of unitary evolutions and projective measurements agree with the
corresponding ordinary complex calculation.
\end{theorem}
\begin{proof}
Each $|\psi_j|^2$ is nonnegative and at most the total sum $1$, so each
component is finite. Apply \cref{prop:st} to the finite norm identity.
The same argument applied to columns of $U$ bounds every entry by $1$ and
shows $U_0^\dagger U_0=I$.

For a standard coordinate vector $e_j$,
\[
\|Pe_j\|^2=P_{jj},\qquad
\|(I-P)e_j\|^2=1-P_{jj}.
\]
Both are nonnegative, so $0\le P_{jj}\le1$ and every entry in the $j$th column
has modulus at most $1$. Reduction preserves the self-adjointness and
idempotence identities. It also preserves the finite sums and products in
\eqref{eq:bornshadow}. Joint probabilities may be computed as squared norms of
products of the relevant projectors and unitaries applied to $\psi$; no
normalization by a possibly infinitesimal intermediate probability is needed.
Applying standard part to this finite product proves the last assertion.
\end{proof}

The conclusion depends on adopting real standard parts as measured
probabilities. It is not a universal theorem about every non-Archimedean
probability interpretation. Conditionalizing on an event whose standard
probability is zero requires additional care, and an infinite or hyperfinite
experiment lies outside the theorem.

A two-level state provides a simple illustration:
\[
\psi=\frac{(1,\varepsilon)}{\sqrt{1+\varepsilon^2}},
\qquad p_2=\frac{\varepsilon^2}{1+\varepsilon^2},\qquad \st(p_2)=0.
\]
Any fixed ordinary finite number of repetitions preserves that zero
standard-part probability for observing the rare outcome. Introducing an
infinite repetition count or an amplifier with infinite resources changes the
operational assumptions; it is not evidence that the original finite
experiment had a different real shadow.

\subsection{What this does and does not say about dynamics}
The theorem associates an ordinary unitary matrix to every chosen finite
surcomplex unitary matrix. It does not prove that the resulting family $U_0(t)$
satisfies a specified Schr\"odinger equation, is strongly continuous, or has a
particular Hamiltonian. Those are differential and topological assertions.
Likewise, an observable can have infinite eigenvalues even though all
normalized-state components are finite. A rule for reading such an observable
cannot be obtained by taking the standard part of an infinite number.

The immediate lesson is favorable but modest: finite quantum algebra is not
obstructed by real-closed non-Archimedean coefficients. To obtain new physics,
one must explain which dynamics or observational rules cannot already be
represented by the ordinary shadow.

\section{Global phases exist, but are not automatically physical time evolution}\label{sec:phases}

\subsection{Finite phases and a global extension}
It would be wrong to dismiss surcomplex quantum theory by saying that no global
surcomplex exponential exists. Ehrlich and Kaplan construct one using an
integer part of the ordered field; for initial surreal structures there is a
canonical choice, the omnific integers $\Oz$ \cite{EK}. The repository's
trigonometry report explicitly includes that construction \cite{RepoTrig}.

For a finite surreal angle $b=b_0+\delta$, with $b_0\in\R$ and
$\delta$ infinitesimal, the usual real phase and an infinitesimal Taylor
expansion give
\[
\operatorname{cis}(b)
 =\ee^{ib_0}\sum_{n\ge0}\frac{(i\delta)^n}{n!}.
\]
The sum is a strong Hahn sum. Every point on the surcomplex unit circle has a
finite angle, for example by the rational half-angle chart together with the
ordinary choice of branch for its finite part. This finite-angle fact does not
assign physical temporal frequencies to arbitrary infinite angles.

Write a general surreal as $b=b_{\rm PI}+b_{\rm fin}$, separating its purely
infinite normal-form terms from its finite terms. With the canonical initial
integer part, the global normalization takes the form
\begin{equation}\label{eq:EKfinite}
\sinEK b=\sin(b_{\rm fin}),\qquad
\cosEK b=\cos(b_{\rm fin}).
\end{equation}
Indeed $b_{\rm PI}/(2\pi)$ belongs to $\Oz$, so the purely infinite part is a
period in the integer-part construction. The corresponding surcomplex
exponential is
\begin{equation}\label{eq:EKexp}
\ExpEK(a+ib)=\exp(a)\bigl(\cosEK b+i\sinEK b\bigr),
\end{equation}
with kernel $2\pi i\Oz$. It extends the ordinary complex exponential and has
the group-homomorphism property. The normalization uses more than the bare
field and local Taylor axioms: it uses the chosen integer part and, for
canonicity here, the surreal initial-subtree structure.

\subsection{An explicit derivative incompatibility}
The global group law is not enough to ensure compatibility with a selected
scalar derivation.

\begin{proposition}[Canonical phase and scalar differentiation]\label{prop:phase}
Let $D$ be a derivation on $\No$ that kills $\R$ and satisfies $D\omega=1$.
The functions in \eqref{eq:EKfinite} do not satisfy the global scalar chain rule
$D(\sinEK x)=\cosEK(x)D(x)$ for every $x\in\No$.
\end{proposition}
\begin{proof}
Since $\omega$ is purely infinite, $\sinEK\omega=0$ and
$\cosEK\omega=1$. Therefore
\[
D(\sinEK\omega)=D(0)=0,
\qquad \cosEK(\omega)D\omega=1.
\]
The asserted global chain rule fails already at $x=\omega$.
\end{proof}

In particular this applies to the normalized Berarducci--Mantova derivation.
It is not a contradiction in either established construction. A derivation of
numbers and the chosen global periodic extension are different structures,
and neither construction claims that these structures can be combined with
all familiar identities without qualification.

There is also a direct real-time warning. For each ordinary real $t$,
$\omega t$ is purely infinite or zero. Hence
\begin{equation}\label{eq:constantphase}
\ExpEK(i\omega t)=1\qquad(t\in\R).
\end{equation}
As a function of an ordinary real time parameter, the right side is constant.
It is therefore not a coefficientwise real-time solution of
$\psi'=i\omega\psi$ with $\psi(0)=1$. A derivative defined using arbitrarily
small \emph{surreal} increments is a different limit operation and does not
repair this real-time substitution. The lesson is to choose a dynamical
function calculus first, rather than infer it from the existence of a phase
homomorphism.

\subsection{Why the naive Taylor series also does not solve the problem}
For an infinite monomial $A$ with $\val(A)<0$, the terms of
\[
\sum_{n\ge0}\frac{(iA)^n}{n!}
\]
have exponents $n\val(A)$ descending without a least element. The series is
not strongly Hahn summable. Even for an ordinary nonzero real angle, the direct
Taylor family is not strongly Hahn summable at exponent zero; its ordinary
complex sum must first be taken in $\C$, as in the finite-phase construction.

An appropriate oscillatory differential algebra, a transferred internal
function space, or a sectorial analytic construction might handle an
infinite-frequency limit. Such a theory may use surcomplex coefficients, but
it must specify which phases are functions, what derivatives act on them, and
how measured finite-time motion is recovered. A global extension of sine alone
is not that theory.

\section{Quantum fields, divergent expansions, and renormalization}\label{sec:qft}

\subsection{The jump from finite matrices to fields is substantial}
Quantum field theory on a curved spacetime needs local field algebras, states,
distributional correlation functions, a suitable spectrum condition, and
renormalized composite observables. In a Hilbert-space presentation it also
needs infinite-dimensional completion and operator domains. Hollands and Wald
provide a precise formulation of these issues for ordinary globally hyperbolic
spacetimes \cite{HW}.

Replacing the coefficient field does not automatically supply a normed-space
completion, countable additivity, a spectral measure, or the hypotheses of a
unitary-evolution theorem. The full fine topology is especially unsuitable for
ordinary countable approximation processes by \cref{prop:discrete}. A chosen
set-sized topology or a coefficientwise construction may be useful, but it is
additional input. Finite Fourier identities and finite Hermitian spectral
algebra do not establish infinite Fourier completeness.

\subsection{Formal perturbative coefficients are a realistic application}
Formal parameters are already a powerful organizing principle in mathematical
physics. Deformation quantization is a well-developed example in which
associative products are built as formal series rather than as convergent
numerical sums \cite{Kontsevich}. A Hahn or surreal realization can place
several such scales in one ordered framework and add fractional powers,
logarithms, or exponential sectors where a justified differential structure is
available.

This can improve the handling of orders of magnitude. It does not make an
arbitrary formal expansion into a unique actual state or scattering amplitude.
For example,
\begin{equation}\label{eq:factorial}
\sum_{n\ge0}(-1)^n n!\,\varepsilon^n
\end{equation}
is a well-defined Hahn series for positive-valuation $\varepsilon$, even though
the associated ordinary power series has zero radius of convergence. The
formal series and an ordinary function with that asymptotic expansion are
different objects.

For positive ordinary $x\to0$, two functions differing by
$C\ee^{-1/x}$ have the same power asymptotic expansion to every order.
A richer transseries can retain the exponentially small sector, but arithmetic
alone does not select its coefficient $C$, a contour prescription, or boundary
conditions. The extension and summation choices in established surreal
integration theories are precisely additional structure, not consequences of
an unadorned infinite sum \cite{CE,BMgerms}.

\subsection{Infinite values do not define products of distributions}
The issue can be seen without a complicated field theory. Let
\[
\delta_\varepsilon(x)=\frac{1}{\varepsilon\sqrt\pi}
                \exp\left(-\frac{x^2}{\varepsilon^2}\right)
\]
for an ordinary positive real regulator $\varepsilon$. Gaussian integration
gives
\begin{equation}\label{eq:deltareg}
\int_\R\delta_\varepsilon(x)\dd x=1,
\qquad
\int_\R\delta_\varepsilon(x)^2\dd x
   =\frac1{\varepsilon\sqrt{2\pi}}.
\end{equation}
The first family tends to the Dirac distribution. The second has a divergent
integral. Naming the right side an infinite surreal value does not construct a
canonical distribution $\delta^2$ or prove regulator independence. The integral
in this example was the ordinary integral for each real regulator; no
unannounced surreal integration rule has been used.

Generalized-function approaches to nonlinear equations do exist, including
frameworks used in general relativity. They retain carefully specified
regularization data and notions of equivalence or association
\cite{SV}. They should not be confused with scalar field extensions. Surreal
coefficients could be incorporated into such a framework, but the mathematical
work lies in defining its operations and its relation to physical
observables.

\subsection{Hawking radiation and information require state evolution}
The Hawking effect can be formulated using ordinary quantum fields on suitable
classical black-hole backgrounds, with the usual care about states and
renormalized observables \cite{HW}. A black-hole information proposal requires
an account of the full evolution and of the observables accessible to relevant
observers. Simply making an interior curvature or an intermediate amplitude
surcomplex neither supplies that evolution nor proves its unitarity.

A successful asymptotic calculation could certainly assist: it might control a
near-horizon mode expansion, a late-time tail, or a nonperturbative correction.
To claim a physical information-resolution mechanism, however, one would still
need to identify a state map, prove the relevant consistency properties, and
relate the result to finite measurement outcomes. Those are the missing
physical statements, not missing names for large numbers.

\section{Existing proposals and comparison with other frameworks}\label{sec:comparison}

\subsection{A direct surreal-singularity proposal deserves a direct test}
Nieto's \emph{Some Mathematical and Physical Remarks on Surreal Numbers}
explicitly discusses gravitational singularities and proposes that the
availability of infinite surreal values removes the usual singularity problem
\cite{Nieto}. In the Schwarzschild discussion, the argument passes from a
larger scalar domain to a claim of disappearance of the singularity, without a
corresponding construction of a regular differential spacetime or a geodesic
continuation.

The calculations in \cref{sec:schwarzschild,sec:infall} identify the missing
step. At a nonzero infinitesimal radius the metric is evaluable, but curvature
is still infinite in the real-boundedness sense and the zero-radius pole is
unchanged. No continuation law follows. Thus the proposed inference is not a
singularity-resolution theorem. This criticism concerns that inference, not
the wider possibility that a future theory using surreal variables could make
nontrivial physical predictions.

\subsection{Hahn fields and transseries often suffice}
For a model using finitely many named scales and set-sized supports, a
suitable Hahn field already provides exact order comparison, ramified powers,
and support-certified algebra. Logarithmic-exponential transseries add an
asymptotic differential language, with established links to surreal functions
\cite{BMgerms}. Full $\No$ offers a universal ambient hierarchy, distinguished
normal forms and simplicity structure, and room to enlarge the scales.

Those are mathematical benefits, but they should be evaluated against a
simpler competitor. If a calculation is unchanged when performed in a small
Hahn field, its success does not demonstrate that the full proper class of
surreals is physically necessary. A positive result should identify whether
the essential contribution is the larger field, a support theorem, a new
extension operator, or simply a more transparent notation for an existing
asymptotic argument.

\subsection{Hyperreal reformulations carry different extra structure}
An ordered non-Archimedean field can sit inside a surreal ambient field, but
this does not identify the accompanying structures. A hyperreal formulation
usually specifies a nonstandard enlargement, internal functions and sets, and
a transfer principle for that enlarged structure. These are more than a field
embedding. The embedding and universality results reviewed by Ehrlich are
valuable precisely because one must state which structure is preserved
\cite{Ehrlich}.

For physics, a nonstandard reformulation may provide a rigorous way to speak
about very fine discretizations and then extract ordinary observables. It does
not automatically change an incomplete classical spacetime into a complete
one. Transferring a classical equation generally preserves its content within
the stated transfer scheme; a new physical resolution requires a changed
model or a changed interpretation of its generalized solutions.

\subsection{Generalized functions and complex continuation solve different problems}
Distributional or generalized-function geometry addresses weak solutions and
nonlinear expressions in fields with limited regularity. Its equivalence
relations, test objects, and product rules are not supplied by $\No$
\cite{SV}. Conversely, a complex contour can detour around a pole, but an
analytic continuation along that contour is not necessarily a real Lorentzian
trajectory. Ordinary $\C$ already permits such detours; changing to $\No[i]$
does not specify which continuation, boundary condition, or causal history is
physical.

There is also no warrant for treating every non-Archimedean theory as the same
proposal. In $p$-adic AdS/CFT, for example, a $p$-adic boundary and a tree-like
bulk are explicitly chosen and correlation functions are computed
\cite{Gubser}. That is a concrete alternative geometric model. It is neither
an ordered extension of the real spacetime metric nor evidence that surreal
numbers remove a Schwarzschild singularity.

\begin{longtable}{@{}>{\raggedright\arraybackslash}p{0.23\textwidth}
  >{\raggedright\arraybackslash}p{0.34\textwidth}
  >{\raggedright\arraybackslash}p{0.35\textwidth}@{}}
\caption{Different extensions answer different mathematical questions. The
last column records a missing implication, not a blanket impossibility.}
\label{tab:comparison}\\
\toprule
Framework & Principal useful structure & What does not follow automatically\\
\midrule
\endfirsthead
\toprule
Framework & Principal useful structure & What does not follow automatically\\
\midrule
\endhead
\bottomrule
\endfoot
Real/complex asymptotics & Limits, error bounds, analytic continuation,
well-developed PDE and operator methods & Existence of a quantum-gravity
completion of a classical singularity\\[5pt]
Hahn or transseries coefficients & Exact hierarchy of scales and
support-controlled formal calculations & Convergence, branch selection, or
physical boundary conditions\\[5pt]
Full surreal/surcomplex scalars & An exceptionally broad ordered/exponential
ambient system and algebraic closure & An ordinary continuum topology,
measurement rule, or globally compatible calculus\\[5pt]
Nonstandard enlargement & Specified transfer and internal structures with
ordinary shadows & New predictions merely from reformulating an old model\\[5pt]
Generalized-function geometry & A specified notion of weak field and
regularized nonlinear operation & Uniqueness or regulator independence without
additional hypotheses\\[5pt]
Modified gravitational dynamics & New field equations or effective sources,
possibly a regular core & Empirical validity, global stability, or unitarity
from local regularity alone\\
\end{longtable}

\section{Physical interpretation, units, and falsifiability}\label{sec:interpretation}

\subsection{Which structure is supposed to be physical?}
A scalar replacement must say whether it changes the set of events, values of
fields on an unchanged spacetime, the codomain of measurement results, or only
an approximation calculus. These choices lead to different theories.
For example, finite surreal field values on a real manifold can have an
ordinary shadow while retaining formal infinitesimal corrections. A manifold
modeled directly on a non-Archimedean field instead requires a new account of
paths, integration and initial data. There is no single object called ``the
surreal spacetime'' determined by the field axioms.

The simplicity relation and the birthday of a surreal number are extra
mathematical structure, not automatically physical age. A rule identifying an
ordinal birthday with time would need to explain how ordinary time translation,
changes of units, and coordinate covariance act on it. No such identification
is needed for the coefficient-based program developed here.

\subsection{Units cannot be ignored when assigning infinitesimal values}
A physical length is represented by a numerical value \emph{and} a unit. The
statements $r/m=\varepsilon$ or $\ell_P/m=\varepsilon^\beta$ are
well-defined scale comparisons. The statement ``the radius is $\omega^{-1}$''
without a unit or reference length is not. Multiplication by a fixed positive
real conversion factor preserves infinitesimality, but a complete physical
model must also specify which dimensional constants and dimensionless ratios
are held fixed in its asymptotic family.

There is no physical reason, from the field alone, to prefer $\omega^{-1}$ to
$\omega^{-2}$ or $\exp(-\omega)$. Their radically different orders become
useful when dictated by an equation or a matching condition. Choosing one as a
fundamental cutoff without such a principle adds unexplained structure.

\subsection{A new scalar ontology may have no new finite predictions}
The standard-part theorem and the finite quantum shadow theorem show how a
large numerical extension can reduce to ordinary predictions for a specified
class of observations. That is not a defect: many valuable mathematical
reformulations are empirically equivalent. It does mean that an ontological
claim---``the true amplitude contains an infinitesimal term''---needs an
operational consequence before it becomes a testable physical distinction.

A proposed distinction should survive changes of the auxiliary embedding,
parameter normalization, and summation prescription whenever those are merely
representational choices. If the prediction changes when an arbitrary
infinitesimal is replaced by another without an accompanying change of
physical inputs, the model has not yet separated physics from its regulator.
If the choice is intended to be physical, it must be promoted to a law or
measurable parameter rather than hidden in notation.

\subsection{The relevant notion of singularity resolution must be stated}
A theory might allow infinite intrinsic quantities but maintain finite detector
responses. Another might replace a spacetime endpoint by a transition between
quantum states. A third might define a weak geometric extension. These are
logically possible strategies, and the present analysis does not reject them
by definition.

Each strategy has a distinct burden of proof. The first must derive detector
responses rather than merely postulate a standard part of a nonexistent finite
value. The second needs a dynamical state map and an interpretation of time.
The third needs a regularity class and an evolution or matching principle.
In all cases the physical gain would be the successful construction and its
consequences, not the permission to write an infinite scalar.

\section{A concrete research program for this repository}\label{sec:program}

\subsection{Recommended architecture}
The most defensible starting point is
\begin{equation}\label{eq:architecture}
\boxed{\begin{gathered}
\text{ordinary spacetime}\\
+\;\text{controlled differential Hahn/transseries coefficients}\\
+\;\text{explicit real-shadow map}
\end{gathered}}
\end{equation}
This preserves the classical geometric domain while enriching the asymptotic
fields. A fully non-Archimedean spacetime can be studied separately, but should
not be required before the coefficient approach has produced a demonstrable
advantage.

The coefficient API should record the value group, support certificates,
coefficient domains, conjugation, real form, and the derivatives under which
the space is closed. It should distinguish a spacetime derivative from a scalar
derivation and a formal parameter derivative at the type or interface level.
The repository's existing finite algebra, Neumann support, rational-jet, and
standard-part layers are useful prerequisites, but their stated current scope
must be respected \cite{RepoRoot}.

\subsection{Benchmark 1: invariant symbolic geometry}
Implement the metric, inverse, connection, curvature, and contractions for
finite-dimensional matrices over a commutative differential coefficient ring.
The first formal target should be \cref{thm:einstein}, since it relies only on
finite algebra, a formal inverse, and commuting derivations. The first explicit
tests should be Schwarzschild in both static and ingoing coordinates and the
Kasner family.

Success means more than producing a large expression. The tool should retain
units, verify agreement of invariant quantities between charts, certify which
denominators are units, and report when a requested evaluation encounters a
zero denominator or unsupported expansion. In particular it should distinguish
$g_{rr}$ at the Schwarzschild horizon from \eqref{eq:Keps} at the central
curvature pole.

\subsection{Benchmark 2: certified dominant balances and matching}
Implement the support test and exact finite remainder in
\cref{prop:matching}. The desired output is a proof-carrying classification of
outer, transition, and inner scales, together with the exact rational model
from which the expansions came. Extend the test to several small parameters
and to fractional-power ansatzes.

A further benchmark is a reduced gravitational or matter ODE known to have a
controlled singular expansion. One should compute its coefficients, prove the
support pattern inductively, and compare with a real solution through a
remainder theorem or a valid summation construction. Formal recurrence alone
is not an existence theorem. In applying a surreal integration result, verify
its nonresonance, function-class, and branch hypotheses explicitly.

\subsection{Benchmark 3: a genuine Einstein--matter construction}
A substantive research target would be to formulate a selected Fuchsian or
quiescent Einstein--matter construction in a differential Hahn setting and
prove that its shadow and remainder bounds reproduce the known classical
solution. The constraints, gauge conditions, and compatibility of evolution
with the constraints must be included. Only then should one ask whether a
larger support class yields a genuinely stronger result.

An alternative target is a controlled inner-layer model with a finite physical
scale. Its source should come from a stated action or effective equations,
not be chosen solely to cancel a curvature pole. One must examine causal
structure, regularity, energy conditions, stability under perturbations, and
matching to the exterior. A finite central scalar is one benchmark among
several, not the complete result.

\subsection{Benchmark 4: compatible quantum dynamics}
Begin with finite systems whose Hamiltonian, real-time evolution and phase
calculus are specified together. Prove reduction to the ordinary Schr\"odinger
equation under common-domain finite-coefficient hypotheses. Treat large-phase
or long-time scalings as separate asymptotic regimes.

The explicit test \eqref{eq:constantphase} should be included in any proposed
phase API: a construction that silently turns $\ExpEK(i\omega t)$ into a
real-time infinite-frequency oscillator is mixing incompatible conventions.
For a field-theoretic extension, specify states, positivity, test functions,
renormalization conditions and the topology of limits before asserting
unitarity or information recovery.

\subsection{What would count as a convincing advance?}
A mathematically convincing advance would be a new controlled solution class,
a sharper singular-asymptotic estimate, a support-based algorithm that detects
invalid expansions, or a rigorous matching theorem that is materially easier
or stronger in the new language. Such results would justify surreal methods
even without any new physical prediction.

A physically convincing advance would additionally identify a well-defined
finite observable, show how it is computed, prove recovery of the appropriate
classical or quantum regime, and explain why the result is independent of
unphysical choices. For a singularity-resolution claim it would specify what
continues through the formerly singular regime and prove the relevant
continuation or dynamical theorem. The present article offers testable
mathematical targets, not a completed quantum-gravity model.

\section{Final assessment}\label{sec:conclusion}

\paragraph{As a language for near-singularity behavior: yes, with substantial potential.}
Surreal and surcomplex methods can record distinct infinite and infinitesimal
scales exactly; fractional powers, logarithms, exponentials and suitably
supported series can expose structure that a single divergent limit suppresses.
The infall and regular-core examples demonstrate concrete applications. The
existing differential and integration theories provide a serious foundation
for selected extensions beyond elementary examples.

\paragraph{As an automatic cure for black-hole singularities: no.}
The Schwarzschild invariant remains $48m^2/r^6$, the proper-time invariant
remains $64/(27u^4)$, and the tidal equations retain their singular coefficients.
A nonzero infinitesimal argument is not the singular endpoint. A regular-core
model changes the metric and its effective source; making its core scale
infinitesimal does not make its central curvature finite in real units.

\paragraph{As a fundamental replacement for all real and complex physics: a
possible research direction, not an established physical theory.}
Finite algebra transfers well. Topology, differentiation, integration,
infinite-dimensional quantum theory, and empirical interpretation require
explicit choices that cannot be inferred from the field alone. There are
already rich positive mathematical constructions, but their compatibility
must be proved rather than assumed. The phase counterexample illustrates this
requirement sharply.

The best-supported next step is therefore to use surreal ideas to strengthen
\emph{the analysis of physical models}, especially multiscale and singular
perturbation problems, while retaining a clear ordinary observational limit.
Such a program could produce valuable mathematics and eventually help a
physical theory of singular regimes. The larger numerical universe supplies
expressive resources; dynamics and evidence must supply the physics.

\clearpage
\appendix
\section{Derivation of the curvature formulas}\label{app:curvature}

For
\[
\dd s^2=-f(r)\dd t^2+\frac{\dd r^2}{f(r)}+r^2\dd\Omega^2,
\]
work first in a region $f>0$ and choose an orthonormal coframe
\[
e^0=\sqrt f\,\dd t,\quad e^1=f^{-1/2}\dd r,\quad
 e^2=r\,\dd\theta,\quad e^3=r\sin\theta\,\dd\phi.
\]
With the curvature convention used in \cref{thm:einstein}, the independent
orthonormal components, up to the Riemann symmetries, may be taken as
\begin{align*}
R_{0101}&=\tfrac12 f'',&
R_{0202}=R_{0303}&=\frac{f'}{2r},\\
R_{1212}=R_{1313}&=-\frac{f'}{2r},&
R_{2323}&=\frac{1-f}{r^2}.
\end{align*}
A simultaneous overall sign change from another Riemann convention does not
change $\Kret$. Each listed squared component has multiplicity four in the
full contraction. Four entries have magnitude $f'/(2r)$, one has magnitude
$f''/2$, and one has magnitude $(1-f)/r^2$. Therefore
\[
\Kret=4\left[\frac{(f'')^2}{4}
          +4\frac{(f')^2}{4r^2}+\frac{(1-f)^2}{r^4}\right],
\]
which is \eqref{eq:Kgeneral}. Contracting before squaring gives
\begin{equation}\label{eq:Rgeneral}
\Scal=-f''-\frac{4f'}r+\frac{2(1-f)}{r^2},
\qquad G^t{}_t=G^r{}_r=\frac{rf'+f-1}{r^2}.
\end{equation}
For $f=1-2M(r)/r$, the second formula becomes
$G^t{}_t=-2M'(r)/r^2$.

These coordinate expressions are rational differential identities in $f$ and
$r$. They continue to regions $f<0$ where the metric is nondegenerate, and to a
regular horizon when a nonsingular coordinate chart such as
\eqref{eq:ef} is used. The static orthonormal coframe itself is not being
asserted to exist with that real square root inside the horizon.

For a Kasner metric, the analogous components have magnitudes
$p_j(p_j-1)/\tau^2$ for the time--space planes and $p_jp_k/\tau^2$ for the
spatial planes. Counting the same multiplicity four gives
\eqref{eq:kasnerK}. The vacuum conditions in \eqref{eq:kasner} follow by
contracting these components.

For Schwarzschild, the time--radial component yields the radial stretching
coefficient $2m/r^3$, and the transverse components yield $-m/r^3$ in the
geodesic-deviation equation. A radial boost preserves these principal tidal
coefficients. This gives \eqref{eq:jacobi} in a radially freely falling frame
rather than only in the static frame.

\clearpage
\section{Audit, provenance, and verification limits}\label{app:audit}

\subsection{Repository inspection}
The repository tree was fetched through the GitHub connector, and the following
files were read at the pinned commit recorded in \cref{sec:repo}:
\begin{center}
\begin{tabular}{@{}l@{}}
\repo{README.md}\\[3pt]
\repo{docs/README.md}\\[3pt]
\repo{docs/surcomplex/analysis/README.md}\\[3pt]
\repo{docs/surcomplex/trigonometry/README.md}
\end{tabular}
\end{center}
The inspection was targeted to the scalar, analytic, topological and
formalization claims relevant to this physics question. It was not a complete
line-by-line audit of every report or a machine verification of their theorems.
The current root README's distinction between generic formalized prerequisites
and the still-missing surreal construction was retained in the assessment.

\subsection{Literature scope}
The source review emphasized original mathematical constructions, original
relativity papers, and technical reviews by researchers in the relevant
fields. It checked a direct published surreal-singularity proposal rather than
assuming that no such proposal exists. The search was targeted, not an
exhaustive systematic literature review; absence of a predictive model in the
inspected sources is not proof that no other proposal has ever been written.

The Ehrlich--Kaplan erratum was checked. It concerns incorrectly numbered
in-text bibliography references in the published version, not withdrawal of
the global phase construction \cite{EKErratum}. The discussion of that
construction here uses the standard $\cos b+i\sin b$ convention, verified by
its value $1$ at $b=0$ and its addition law.

\subsection{Status of the results in this article}
The standard-part, pole, formal-reduction, support, spectral, shadow and
phase-compatibility arguments are proved in the text. They are supplied as
explicit mathematical deductions for this assessment; no priority claim is
made. The physical equations used in the examples are standard models, not
newly discovered black-hole solutions. The research architecture and benchmarks
are proposals, not implemented or experimentally validated theories.

The accompanying \texttt{verify.py} uses SymPy to check finite symbolic
identities: the Schwarzschild curvature and Ricci scalar, the proper-time
curvature coefficient, the tidal indicial roots, Kasner special cases, the
regular-core curvature and source, the boundary-layer derivative, the matching
identities, and the finite geometric remainders. It also independently builds
the static metric's Christoffel symbols, Riemann tensor and Ricci tensor, and
checks both curvature contractions. All 45 finite checks passed in the
recorded run. These checks are useful against algebraic transcription errors; they
are not proofs of the general topological, formal-series, or physical claims.
No Lean formalization of this article is included or claimed.

\clearpage
\section{Notation and assumptions at a glance}\label{app:notation}
\begin{longtable}{@{}>{\raggedright\arraybackslash}p{0.23\textwidth}
 >{\raggedright\arraybackslash}p{0.71\textwidth}@{}}
\toprule
Symbol or phrase & Meaning in this article\\
\midrule
\endfirsthead
\toprule Symbol or phrase & Meaning in this article\\\midrule
\endhead
\bottomrule
\endfoot
$\No$, $\No[i]$ & Full proper-class surreal field and its complexification.
Not an ordinary set-sized metric space.\\[4pt]
$F_\Gamma$, $K_\Gamma$ & Set-sized real and complex Hahn fields with well-ordered
support and small monomials $t^\gamma$ for $\gamma>0$.\\[4pt]
$\val$ & Least exponent of a nonzero Hahn series; $\val(0)=+\infty$ by
convention when needed.\\[4pt]
$\OO_F$, $\mm_F$ & Elements bounded by an ordinary real bound, and elements
smaller in magnitude than every positive ordinary real.\\[4pt]
$\st$ & Real or complex standard part, defined only on finite elements.\\[4pt]
$\varepsilon$ & A positive dimensionless infinitesimal when working in a Hahn
field; a positive real regulator when explicitly discussing a real family.\\[4pt]
$h$, $\pi_0$ & Formal parameter fixed by coordinate derivatives, and reduction
modulo $h$ in the formal Einstein theorem.\\[4pt]
$\partial_\mu$, $D$ & Coefficientwise ordinary spacetime derivative, and a
specified scalar derivation on surreal numbers. They are not interchangeable.\\[4pt]
$m$, $\ell$, $\ell_P$ & Geometrized mass $GM/c^2$, core length scale, and
Planck length.\\[4pt]
$\tau$, $u$ & Proper time in length units and remaining proper time
$u=\tau_*-\tau$.\\[4pt]
$\Kret$ & Kretschmann scalar $R_{abcd}R^{abcd}$; not the coefficient field $K$.\\[4pt]
$\Oz$, $\ExpEK$ & Canonical initial integer part of $\No$, and the
integer-part-based global surcomplex exponential.\\[4pt]
Formal solution & An identity coefficient by coefficient. It does not by itself
assert convergence, summability to a real solution, or physical validity.\\
\end{longtable}

\clearpage
\begin{thebibliography}{99}
\raggedright
\addcontentsline{toc}{section}{References}

\bibitem{RepoRoot}
V.~Reshetnikov, \emph{Surreal}, repository root README, pinned snapshot
identified in \cref{sec:repo}.
\repo{README.md}. Accessed 21 September 2026.

\bibitem{RepoDocs}
V.~Reshetnikov, \emph{The surreal and surcomplex reports}, repository catalogue,
same snapshot. \repo{docs/README.md}. AI-assisted research documentation;
accessed 21 September 2026.

\bibitem{RepoAnalysis}
V.~Reshetnikov, \emph{Surcomplex analysis: holomorphic functions on
$K=\No[i]$}, report README, same snapshot.
\repo{docs/surcomplex/analysis/README.md}. Accessed 21 September 2026.

\bibitem{RepoTrig}
V.~Reshetnikov, \emph{Trigonometry on the Surcomplex Plane}, report README,
same snapshot. \repo{docs/surcomplex/trigonometry/README.md}.
Accessed 21 September 2026.

\bibitem{BM}
A.~Berarducci and V.~Mantova, Surreal numbers, derivations and transseries,
\emph{Journal of the European Mathematical Society} \textbf{20} (2018),
339--390. \doi{10.4171/JEMS/769}; \arxiv{1503.00315}.

\bibitem{BMgerms}
A.~Berarducci and V.~Mantova, Transseries as germs of surreal functions,
\emph{Transactions of the American Mathematical Society} \textbf{371} (2019),
3549--3592. \doi{10.1090/tran/7428}; \arxiv{1703.01995}.

\bibitem{EK}
P.~Ehrlich and E.~Kaplan, Surreal ordered exponential fields,
\emph{Journal of Symbolic Logic} \textbf{86} (2021), 1066--1115.
\doi{10.1017/jsl.2021.59}; \arxiv{2002.07739}.
In particular, Section 11 on trigonometric fields and surcomplex exponentiation.

\bibitem{EKErratum}
P.~Ehrlich and E.~Kaplan, Surreal ordered exponential fields---Erratum,
\emph{Journal of Symbolic Logic} \textbf{87} (2022), 871.
\doi{10.1017/jsl.2022.5}.

\bibitem{Ehrlich}
P.~Ehrlich, The absolute arithmetic continuum and the unification of all
numbers great and small, \emph{Bulletin of Symbolic Logic} \textbf{18} (2012),
1--45. \doi{10.2178/bsl/1327328438}.

\bibitem{CE}
O.~Costin and P.~Ehrlich, Integration on the surreals,
\emph{Advances in Mathematics} \textbf{452} (2024), article 109823.
\arxiv{2208.14331}, version 5 consulted.

\bibitem{Penrose}
R.~Penrose, Gravitational collapse and space-time singularities,
\emph{Physical Review Letters} \textbf{14} (1965), 57--59.
\doi{10.1103/PhysRevLett.14.57}.

\bibitem{Witten}
E.~Witten, Light rays, singularities, and all that,
\emph{Reviews of Modern Physics} \textbf{92} (2020), 045004.
\doi{10.1103/RevModPhys.92.045004}; \arxiv{1901.03928}.

\bibitem{Raychaudhuri}
A.~Raychaudhuri, Relativistic cosmology. I,
\emph{Physical Review} \textbf{98} (1955), 1123--1126.
\doi{10.1103/PhysRev.98.1123}.

\bibitem{Nieto}
J.~A.~Nieto, Some mathematical and physical remarks on surreal numbers,
\emph{Journal of Modern Physics} (2016).
\doi{10.4236/jmp.2016.715188}; \arxiv{1611.09699}.
The black-hole discussion appears near the end of the 19-page preprint.

\bibitem{Hayward}
S.~A.~Hayward, Formation and evaporation of non-singular black holes,
\emph{Physical Review Letters} \textbf{96} (2006), 031103.
\doi{10.1103/PhysRevLett.96.031103}; \arxiv{gr-qc/0506126}.

\bibitem{CR}
R.~Carballo-Rubio, F.~Di Filippo, S.~Liberati and M.~Visser,
Geodesically complete black holes,
\emph{Physical Review D} \textbf{101} (2020), 084047.
\doi{10.1103/PhysRevD.101.084047}; \arxiv{1911.11200}.

\bibitem{Donoghue}
J.~F.~Donoghue, General relativity as an effective field theory: The leading
quantum corrections, \emph{Physical Review D} \textbf{50} (1994), 3874--3888.
\doi{10.1103/PhysRevD.50.3874}; \arxiv{gr-qc/9405057}.

\bibitem{DonoghueReview}
J.~F.~Donoghue, Quantum general relativity and effective field theory,
\arxiv{2211.09902}. Review contribution to the \emph{Handbook of Quantum Gravity}.

\bibitem{AR}
L.~Andersson and A.~D.~Rendall, Quiescent cosmological singularities,
\emph{Communications in Mathematical Physics} \textbf{218} (2001), 479--511.
\doi{10.1007/s002200100406}; \arxiv{gr-qc/0001047}.

\bibitem{Dafermos}
M.~Dafermos, The interior of charged black holes and the problem of uniqueness
in general relativity, \arxiv{gr-qc/0307013}.

\bibitem{DL}
M.~Dafermos and J.~Luk, The interior of dynamical vacuum black holes I:
The $C^0$-stability of the Kerr Cauchy horizon, \arxiv{1710.01722}.
The use here is restricted to the hypotheses and interior conclusion of this work.

\bibitem{HW}
S.~Hollands and R.~M.~Wald, Quantum fields in curved spacetime,
\emph{Physics Reports} (2015). \doi{10.1016/j.physrep.2015.02.001};
\arxiv{1401.2026}.

\bibitem{Kontsevich}
M.~Kontsevich, Deformation quantization of Poisson manifolds, I,
\arxiv{q-alg/9709040}.

\bibitem{SV}
R.~Steinbauer and J.~A.~Vickers, The use of generalised functions and
distributions in general relativity,
\emph{Classical and Quantum Gravity} \textbf{23} (2006), R91--R114.
\arxiv{gr-qc/0603078}.

\bibitem{Gubser}
S.~S.~Gubser, J.~Knaute, S.~Parikh, A.~Samberg and P.~Witaszczyk,
$p$-adic AdS/CFT, \emph{Communications in Mathematical Physics} (2017).
\doi{10.1007/s00220-016-2813-6}; \arxiv{1605.01061}.

\end{thebibliography}
\end{document}
